% !TEX program = lualatex
% Generated by new_lecture_note.sh
\documentclass[11pt]{article}

\usepackage{fontspec}
\usepackage{microtype}
\usepackage{mathtools,amssymb,amsfonts,amsthm,bm}
\usepackage{enumitem}
\usepackage{booktabs,array,xcolor}
\usepackage{geometry,fancyhdr,setspace}
\usepackage[hidelinks]{hyperref}
\usepackage[nameinlink,noabbrev]{cleveref}

\geometry{margin=1in,headheight=15pt}
\setstretch{1.12}
\setlength{\parindent}{0pt}
\setlength{\parskip}{0.55em plus 0.12em minus 0.08em}
\setlength{\abovedisplayskip}{0.8em plus 0.2em minus 0.1em}
\setlength{\belowdisplayskip}{0.8em plus 0.2em minus 0.1em}
\allowdisplaybreaks[3]
\setlist{leftmargin=*,topsep=0.35em,itemsep=0.15em,parsep=0pt}

% ---------- Metadata ----------
\newcommand{\studentname}{Keshav Ramamurthy}
\newcommand{\coursename}{Math 118}
\newcommand{\lecturenumber}{01}
\newcommand{\lecturetitle}{Lecture 01}
\newcommand{\lecturedate}{\today}

% ---------- Reusable notation ----------
\newcommand{\N}{\mathbb{N}}
\newcommand{\Z}{\mathbb{Z}}
\newcommand{\Q}{\mathbb{Q}}
\newcommand{\R}{\mathbb{R}}
\newcommand{\C}{\mathbb{C}}
\newcommand{\F}{\mathbb{F}}
\newcommand{\K}{\mathbb{K}}
\newcommand{\E}{\mathbb{E}}
\newcommand{\Prob}{\mathbb{P}}
\newcommand{\1}{\mathbf{1}}
\newcommand{\eps}{\varepsilon}
\newcommand{\dd}{\,\mathrm{d}}
\newcommand{\abs}[1]{\left\lvert#1\right\rvert}
\newcommand{\norm}[1]{\left\lVert#1\right\rVert}
\newcommand{\ip}[2]{\left\langle#1,#2\right\rangle}
\newcommand{\set}[1]{\left\{#1\right\}}
\newcommand{\given}{\,\middle\vert\,}
\newcommand{\indep}{\perp\!\!\!\perp}
\newcommand{\lto}{\longrightarrow}
\newcommand{\deriv}[2]{\frac{\mathrm{d}#1}{\mathrm{d}#2}}
\newcommand{\pderiv}[2]{\frac{\partial#1}{\partial#2}}
\DeclareMathOperator{\Aut}{Aut}
\DeclareMathOperator{\Hom}{Hom}
\DeclareMathOperator{\Ker}{ker}
\DeclareMathOperator{\im}{im}
\DeclareMathOperator{\rank}{rank}
\DeclareMathOperator{\nullity}{nullity}
\DeclareMathOperator{\Span}{span}
\DeclareMathOperator{\Var}{Var}
\DeclareMathOperator{\Cov}{Cov}

% ---------- Note environments ----------
\newtheoremstyle{notestyle}{0.7em}{0.7em}{\itshape}{}{}{\bfseries}{.5em}{}
\theoremstyle{notestyle}
\newtheorem{theorem}{Theorem}[section]
\newtheorem{lemma}[theorem]{Lemma}
\newtheorem{proposition}[theorem]{Proposition}
\newtheorem{corollary}[theorem]{Corollary}
\theoremstyle{definition}
\newtheorem{definition}[theorem]{Definition}
\newtheorem{example}[theorem]{Example}
\newtheorem{remark}[theorem]{Remark}
\newenvironment{claim}{\par\noindent\textbf{Claim.}\ }{\par}
\newenvironment{idea}{\par\noindent\textit{Idea.}\ }{\par}
\newenvironment{proofsketch}{\par\noindent\textbf{Proof sketch.}\ }{\par}

% ---------- Header ----------
\pagestyle{fancy}
\fancyhf{}
\fancyhead[L]{\coursename}
\fancyhead[C]{\lecturetitle}
\fancyhead[R]{\studentname}
\fancyfoot[C]{\thepage}

\title{\vspace{-1.5em}\textbf{\coursename -- \lecturetitle}}
\author{\studentname}
\date{\lecturedate}

\begin{document}
\maketitle

\section{Main idea}
\textit{Write the one-sentence point of today's lecture here.}
Basic information about everything.
\par Chapter 0 is preliminaries, inner product spaces
\section{Definitions and notation}
\begin{definition}[Previous Notes]
    Recall in $\mathbb{C}$, $$z = r e^{i \theta} = r \left(\cos\theta + i \sin \theta\right)$$
    Notice this makes multiplication pretty nice in $\mathbb{C}$, as angles combine
    easily

% Write the definition and one sentence explaining why it matters.
\end{definition}
\begin{definition}[Power Series of e (Taylor Series)]
\label{def:label}
For $z \in \mathbb{C}$, $e^z$ is defined by $$e^z = \sum_{n = 0}^{\infty}
\frac{z^n}{n!}$$ 
$$\sum_{n=0}^{\infty} \frac{(z+w)^n}{n!} = \sum_{n=0}^{\infty}\frac{1}{n!}\left(\sum_{k=0}^{n}\binom{n}{k}z^kw^{n-k}\right)$$
The $n!$ cancel out and then you can rearrange into(NUMBERING THE LATTICE)
$$\sum_{k=0}^{\infty}\sum_{n=k}^{\infty} \frac{1}{k!(n-k)! z^kw^{n-k}}$$
Rearrange $$j = n-k \rightarrow \sum_{k=0}^{\infty}\sum_{j=0}^{\infty}\frac{1}{k!}z^k
\frac{1}{j!}{w^{j}} = \left(\sum_{k=0}^{\infty}\frac{1}{k!}z^k\right)\left(\sum_{j=0}^{\infty}\frac{1}{j!}w^{j}\right)$$
\end{definition}
\begin{definition}[Derivation of De Moivre's in a different way]
\label{def:label}
If you look at $$e^{i\theta}=\sum_{n=0}^{\infty} \frac{(i \theta)^n}{n!} $$
which is 
$$\left(1 - \frac{\theta^2}{2!} + \frac{\theta^{4}}{4!} + ...\right) + i \left(\theta -
\frac{\theta^3}{3!}+\frac{\theta^5}{5!} - ...\right) = \cos \theta + i \sin \theta$$
\end{definition}
\begin{definition}[Remember These Properties of Conjugation]
\label{def:label}
$$\overline{zw} = \overline{z}\cdot \overline{w}, \overline{\left(\frac{z}{w}\right)} =
\frac{\overline{z}}{\overline{w}}$$
If $$f(z) = \sum_{n=0}^{\infty} c_n (z-a)^n, c_n a \in \mathbb{R} $$
then $$f(\overline{z}) = \overline{f(z)}$$
\end{definition}

\begin{definition}[Vector Space]
\label{def:label}
Vector space over a field like real or complex where you can perform addition and
multiply by scalars, etc...
\par $\mathbb{C}^{n}$ is basically a vector space over $\mathbb{R}$ of dimension 2n. (but
only real operations)
\par $\textbf{Infinite Dimensional Spaces}$ (like the set of all polynomials)
\par For example, $l^2(\mathbb{N}) = \text{space of infinite sequences } (x_1, x_2, \cdots)$
such that the sum $\sum_{n=1}^{\infty}|x_n|^2$ converges.
\par For example, $$l^1(\mathbb{Z}) = \text{doubly infinite, summable sequences}
\sum_{x=1}^{\infty} |x_i| < \infty, \text{ Show that } l^1(\mathbb{N}) \text{ is a
subset of } l^2(\mathbb{N})$$
\par \textbf{Proof: }
If you look at some $x \in l^1(\mathbb{N})$, we have $\sum_{n=1}^{\infty} |x_i| $
\par You can just look at the expansion of $C^2 =
\sum_{a=1}^{\infty}|x_a|\sum_{b=1}^{\infty}|x_b|$ which is just the terms we want plus a lo t
of extraneous absolute value terms bounding $\sum_{i=1}^{\infty}|x_i^2|$
converges.
\par For example, $$L'([a, b]) = f : [a, b] \rightarrow \mathbb{C}\text{ such that } {\int_{a}^{b}|f(x)|dx
\text {converges }}$$ 
\par One can basically show $L^2([a, b]) \le L'([a, b])$ like $f(x) =
\frac{1}{\sqrt{x}}$ which is in $L'([0, 1])$ but not $L^2([0, 1])$
\par For example, $L^4(\mathbb{R}^2)$ is a function $\mathbb{R}^2 \rightarrow
\mathbb{C}$ such that $$\int_{-\infty}^{\infty}\infty_{-\infty}^{\infty}|f(x,y)|^4dx
dy \text { converges} $$ 
\end{definition}

\section{Claims, theorems, and proof sketches}
\begin{claim}
State the claim in one precise sentence.
\end{claim}

\begin{proofsketch}
Record the key idea, the first useful equation, and the step that still needs checking.
\end{proofsketch}

\section{Examples and calculations}
\begin{example}[Lecture example]
Write the setup, the calculation, and the conclusion.
\end{example}

\section{Questions and follow-up}
\begin{itemize}
  \item What is still unclear?
  \item Which exercise or reading should I revisit?
  \item TODO: 
\end{itemize}

\end{document}
